% Options for packages loaded elsewhere
\PassOptionsToPackage{unicode}{hyperref}
\PassOptionsToPackage{hyphens}{url}
%
\documentclass[
  10pt,
]{article}
\usepackage{lmodern}
\usepackage{amssymb,amsmath}
\usepackage{ifxetex,ifluatex}
\ifnum 0\ifxetex 1\fi\ifluatex 1\fi=0 % if pdftex
  \usepackage[T1]{fontenc}
  \usepackage[utf8]{inputenc}
  \usepackage{textcomp} % provide euro and other symbols
\else % if luatex or xetex
  \usepackage{unicode-math}
  \defaultfontfeatures{Scale=MatchLowercase}
  \defaultfontfeatures[\rmfamily]{Ligatures=TeX,Scale=1}
\fi
% Use upquote if available, for straight quotes in verbatim environments
\IfFileExists{upquote.sty}{\usepackage{upquote}}{}
\IfFileExists{microtype.sty}{% use microtype if available
  \usepackage[]{microtype}
  \UseMicrotypeSet[protrusion]{basicmath} % disable protrusion for tt fonts
}{}
\makeatletter
\@ifundefined{KOMAClassName}{% if non-KOMA class
  \IfFileExists{parskip.sty}{%
    \usepackage{parskip}
  }{% else
    \setlength{\parindent}{0pt}
    \setlength{\parskip}{6pt plus 2pt minus 1pt}}
}{% if KOMA class
  \KOMAoptions{parskip=half}}
\makeatother
\usepackage{xcolor}
\IfFileExists{xurl.sty}{\usepackage{xurl}}{} % add URL line breaks if available
\IfFileExists{bookmark.sty}{\usepackage{bookmark}}{\usepackage{hyperref}}
\hypersetup{
  hidelinks,
  pdfcreator={LaTeX via pandoc}}
\urlstyle{same} % disable monospaced font for URLs
\usepackage[a4paper,margin=18mm]{geometry}
\setlength{\emergencystretch}{3em} % prevent overfull lines
\providecommand{\tightlist}{%
  \setlength{\itemsep}{0pt}\setlength{\parskip}{0pt}}
\setcounter{secnumdepth}{-\maxdimen} % remove section numbering

\usepackage{esint}
\usepackage{xeCJK}
\setmainfont{DejaVu Serif}
\setmonofont{DejaVu Sans Mono}
\setCJKmainfont{Noto Serif CJK SC}
\setCJKmonofont{Noto Sans Mono CJK SC}
\setmathfont{Asana Math}
\setmathfont{Latin Modern Math}[range=\lrcorner]
\setlength{\emergencystretch}{3em}
\allowdisplaybreaks
\hypersetup{hypertexnames=false}

\author{}
\date{}

\begin{document}

\hypertarget{signed-characteristic-mass-and-the-missing-bulk-selection}{%
\section{Signed characteristic mass and the missing bulk
selection}\label{signed-characteristic-mass-and-the-missing-bulk-selection}}

28 September 2026. Continuation from \texttt{0fb009c}. The first-order
target is still the infimum over all metrics on the original balanced
family. This text identifies exactly what I55's measurable bulk
selection requires, supplies it in a classical Kähler setting, and tests
its necessity with actual smooth metrics. In particular, a competitive
energy bound is not enough to impose that selection, even after
constant-scalar normalization. This does not produce a coefficient below
four.

\hypertarget{the-signed-measure-behind-i55}{%
\subsection{1. The signed measure behind
I55}\label{the-signed-measure-behind-i55}}

Use the complex orientation on

\[
 M=X\#k\overline{\mathbb{CP}}^{\,2}\#\ell(S^1\times S^3),
 \qquad A=c_1^2(X),\quad Q=k+4\ell,\quad h=Q/A,
\]

where X is one of the complex-hyperbolic covers already used in the
actual family. All curvature norms have the conventions of I54--I55. For
any smooth metric put

\[
 f_g=\frac{s^2/12+4|W^+|^2-|\operatorname{Ric}^{\circ}|^2}{8\pi^2},
 \quad N(g)=\int_M(f_g)_-,\quad T_U=\int_U f_g.
 \tag{1.1}
\]

Chern--Weil gives \(T_M=A-Q\). N is the negative mass of this particular
Levi-Civita characteristic density, not an ADM mass, a characteristic
number by itself, or the torsion energy denoted N in earlier cell
formulas. Let \(R=\int|\operatorname{Ric}|^2\),
\(P=\int|\operatorname{Ric}^{\circ}|^2\), and \(L_s=\int(s-\bar s)^2\).
Then

\[
 \mathcal E=\frac{4P+L_s}{R},\qquad 8\pi^2N\le P.
 \tag{1.2}
\]

Consider h tending to zero and competitive actual metrics
\(\mathcal E=O(h)\). The classical Ricci lower bound in this family is
\(R\ge8\pi^2(A+Q)\). The signature formula gives

\[
 \int|W^+|^2\ge12\pi^2\tau
   =4\pi^2A/3-12\pi^2k.
\]

Combining it with \(R=4P-12\int|W^+|^2+24\pi^2(A-Q)\) and
\(4P\le\mathcal E R\) proves

\[
 R=8\pi^2A+O(Q),\qquad P+L_s=O(Q),\qquad N=O(Q).
 \tag{1.3}
\]

These are bounds on actual metrics, not an optimization of formal
curvature moments. They use the classical lower bound already supplied
in I44.

\textbf{Proposition 1.} In this competitive regime the following
conditions are equivalent:

\[
 \begin{aligned}
 &\text{there are measurable }U_j\text{ with }
 T_{U_j}=A_j+o(Q_j),\quad R_{U_j}=8\pi^2A_j+o(A_j);\\
 &N(g_j)\ge Q_j-o(Q_j).
 \end{aligned} \tag{1.4}
\]

The first condition is precisely the intrinsic sufficient interface in
\href{20260928_localization_and_coupled_filling.md}{I55, Section 4.1}.

\emph{Proof.} Necessity follows from \(T_{M\setminus U}=-Q+o(Q)\). For
sufficiency choose a measurable \(C\subset\{f<0\}\) with negative mass
\(\min(N,Q)\). Such a set exists because the finite measure \((f)_-dV\)
has no atoms. On \(\{f<0\}\),
\(s^2/12\le|\operatorname{Ric}^{\circ}|^2\), so
\(|\operatorname{Ric}|^2\le4|\operatorname{Ric}^{\circ}|^2\).
Consequently \(R_C=O(Q)\). For \(U=M\setminus C\), (1.3) now gives both
required estimates. QED.

The coefficient consequence can also be seen immediately from (1.2):
(1.4) implies \(\liminf\mathcal E/h\ge4\). I55 adds its equality
information. Neither argument proves (1.4) from competitiveness.

\hypertarget{a-genuine-supplier-from-holomorphic-morse-inequalities}{%
\subsection{2. A genuine supplier from holomorphic Morse
inequalities}\label{a-genuine-supplier-from-holomorphic-morse-inequalities}}

For this section let \(M=\operatorname{Bl}_k X\) with its specified
complex structure, with no handles, and X a closed complex-hyperbolic
surface. For every smooth Hermitian metric on the canonical line bundle
write \(\Theta=(i/2\pi)F_{K_M}\). The weak holomorphic Morse inequality
gives

\[
 \limsup_{m\to\infty}\frac{2h^1(M,K_M^m)}{m^2}
 \le\int_{M(\Theta,1)}-\Theta^2. \tag{2.1}
\]

Here \(M(\Theta,1)\) is the locus with one positive and one negative
curvature eigenvalue. The classical input, including its signature
convention, is \href{https://arxiv.org/html/1003.5067}{Demailly,
equations (1.11)--(1.12)}. Only the inequality is needed, not a converse
or a conjectured equality.

In the present case the left side is exactly k. Indeed K\_X is ample,
Kodaira vanishing and Riemann--Roch give

\[
 h^0(X,K_X^m)=\frac{A}{2}m(m-1)+\chi(\mathcal O_X),\qquad m\ge2.
\]

Pluricanonical sections and \(\chi(\mathcal O)\) are unchanged by
blowup, while \(K_M^2=A-k\). Serre duality gives \(h^2(M,K_M^m)=0\): an
antipluricanonical effective divisor would have negative pairing with
the pullback of the ample K\_X. Thus

\[
 h^1(M,K_M^m)=\frac{k}{2}m(m-1),\qquad
 \int_{M(\Theta,1)}-\Theta^2\ge k. \tag{2.2}
\]

If g is Kähler, its induced canonical curvature is minus the Ricci form
divided by \(2\pi\). Squaring removes that sign. The Kähler identity
\(|W^+|^2=s^2/24\) gives

\[
 \Theta^2=\frac{s^2/4-|\operatorname{Ric}^{\circ}|^2}{8\pi^2}\,dV
          =f_g\,dV.
\]

The signature-one set is exactly the negative-density set. Hence

\[
 \boxed{N(g)\ge k\quad\text{for every Kähler metric on }\operatorname{Bl}_k X.}
 \tag{2.3}
\]

This is a consequence of classical holomorphic Morse theory, not a new
Morse theorem. It supplies spatially signed curvature information, in
addition to the familiar integrated Ricci lower bounds. Applying
Proposition 1 proves the I55 selection for competitive Kähler sequences
in this pure-blowup class. Handles and non-Kähler metrics are not
covered. For a conformally Kähler metric the canonical Chern curvature
still satisfies (2.2), but it is no longer the density (1.1). The
distinction is essential in the following construction.

\hypertarget{an-exact-radial-primitive-and-an-actual-loss-of-negative-mass}{%
\subsection{3. An exact radial primitive and an actual loss of negative
mass}\label{an-exact-radial-primitive-and-an-actual-loss-of-negative-mass}}

Take the ordinary one-bolt metric K\_F of I54, with
\(r\ge r_0=\sqrt\mu\), volume element \(2\pi^2r^3dr\), \(F(r_0)=0\),
\(F'(r_0)=2/r_0\), and a full CH outer collar. For \(G=e^{2\phi}K_F\),
define

\[
 q=2F+\frac{rF'}2+rF\phi'-2,\qquad
 B=q^2+(2+q)r^2F(\phi')^2.
 \tag{3.1}
\]

Direct substitution of the conformal Ricci eigenvalues gives the exact
identity

\[
 \boxed{\int_{\{r\in[a,b]\}} f_G\,dV_G=B(b)-B(a).} \tag{3.2}
\]

For a reproducible derivation, let \(r_v=-F''/2-5F'/(2r)\),
\(r_h=-F'/r+4(1-F)/r^2\). The three scaled Ricci eigenvalues, the
horizontal one having multiplicity two, are

\[
\begin{aligned}
 A_0&=r_v-3F\phi''-(2F'+3F/r)\phi',\\
 A_3&=r_v-F\phi''-(2F'+5F/r)\phi'-2F(\phi')^2,\\
 A_H&=r_h-F\phi''-(F'+5F/r)\phi'-2F(\phi')^2.
\end{aligned}
\]

Put \(S=A_0+A_3+2A_H\). Conformal Weyl invariance and
\(s_{K_F}=2(r_v+r_h)\) reduce the density to

\[
 B'=\frac{r^3}{4}\left\{\frac{S^2}{12}
 +\frac{[2(r_v+r_h)]^2}{6}
 -\left(A_0^2+A_3^2+2A_H^2-\frac{S^2}{4}\right)\right\}.
 \tag{3.3}
\]

Expansion proves (3.2); the checking program independently verifies the
identity for arbitrary smooth F and phi. At an ordinary smooth bolt B=1;
on a CH collar B=r\^{}4/4. Thus the relative characteristic integral is
correct at both ends, including its normalization.

For the uncut constant-scalar recovery metric use

\[
 F_\mu=1+\frac{r^2}{6}-\frac{c}{2r^2}+\frac{\mu^3}{3r^4},
 \qquad c=\mu(\mu+2).
\]

It has s=-4 and \(q_0=(r^2-c/r^2)/2\). The function \(B_0=q_0^2\)
decreases from one to zero at \(r_*=c^{1/4}\), then increases. Its
negative mass is exactly one. Equivalently \(8\pi^2f=4(1-c^2/r^8)\), and
direct integration on \([\sqrt\mu,r_*]\) gives one.

Fix a small positive mu. Choose smooth eta supported near r\_* with
\(\eta'(r_*)\ne0\), and set \(\phi=\varepsilon\eta\). The nondegenerate
minimum persists. Its position first shifts so as to cancel the linear
change in q; the remaining value is

\[
 \min B_\varepsilon
 =2r_*^2F_\mu(r_*)\eta'(r_*)^2\varepsilon^2+O_\mu(\varepsilon^3)>0.
 \tag{3.4}
\]

Cut F\_mu to CH in a fixed annulus outside r\_* where B\_0' is strictly
positive. For sufficiently small mu that cutoff remains C2-close to CH
on the annulus, so its density is positive there. Small epsilon
preserves the single inner minimum and creates no additional negative
density outside. The complete smooth cell therefore satisfies

\[
 N(G_\varepsilon)=1-\min B_\varepsilon<1. \tag{3.5}
\]

Its cost J is still strictly above \(8\pi^2\), by I54's complete
conformal-Kähler estimate. It is not a cheaper cell.

Holding mu and epsilon fixed, repetition in covers with k/A tending to
zero gives actual closed metrics with \(\mathcal E=O(k/A)\) but
\(N/k<1\) by a fixed amount. The exact repetition formula is I54 (2.1).
Thus competitiveness alone does not supply Proposition 1, already in the
pure-blowup family.

The fixed-mu qualification cannot be dropped. If the perturbation is
confined to \(r\asymp\mu^{1/4}\) while mu tends to zero, B\_0 at the
support boundaries is O(mu). The accumulated negative variation is then
already at least \(1-O(\mu)\). Formula (3.4) is not uniform for a fixed
perturbation amplitude in that limit; new shoulder minima can appear. In
particular this argument does not refute selection for shrinking
near-minimizers.

\hypertarget{constant-scalar-cells-and-the-original-balanced-family}{%
\subsection{4. Constant-scalar cells and the original balanced
family}\label{constant-scalar-cells-and-the-original-balanced-family}}

The conformal bump in Section 3 disappears under Yamabe minimization of
its original conformal class. A separate construction is needed to test
the scalar-minus-four formulation of the selection target.

\hypertarget{a-scalar-minus-four-cell-with-n-less-than-one}{%
\subsubsection{4.1 A scalar-minus-four cell with N less than
one}\label{a-scalar-minus-four-cell-with-n-less-than-one}}

On a fixed truncation of K\_\{F\_mu\}, solve the radial equation

\[
 (r^3F_\mu u_a')'=\frac23r^3(u_a^3-u_a),\qquad
 u_a(r_0)=1+a,\qquad (r^3F_\mu u_a')(r_0)=0. \tag{4.1}
\]

It is the conformal scalar equation for \(u_a^2K_{F_\mu}\) to have s=-4.
Smooth bolt regularity gives a smooth local solution and smooth
dependence on a; for fixed mu and fixed finite outer radius it stays
positive for sufficiently small a. For a\textgreater0 it initially
exceeds one, and the integrated equation shows that it increases, so
\(u_a'(r_*)>0\). Differentiating at a=0 gives \(v(r_0)=1\) and

\[
 (r^3F_\mu v')'=\frac43r^3v,\qquad v'(r_*)>0.
\]

Use \(\phi_a=\log u_a\) in (3.1). The same minimum calculation now gives

\[
 N_{\rm inner}=1-2r_*^2F_\mu(r_*)v'(r_*)^2a^2+O_\mu(a^3)<1.
 \tag{4.2}
\]

This inner metric is not CH on its outer collar. To complete it without
changing its inner region or its scalar curvature, apply
\href{https://arxiv.org/html/1003.5146}{Delay, Theorem 1.1 and
Propositions 2.2--2.5} on a fixed annulus. The theorem glues
sufficiently close smooth metrics with a smooth supported correction,
interpolating their scalar curvatures, provided the reference adjoint
scalar operator has no kernel.

That nondegeneracy holds for CH. In geometric Laplacian convention a
kernel function satisfies \(\operatorname{Hess}f-(\Delta f)g+fg=0\),
hence \(\operatorname{Hess}f=(f/3)g\). Commuting derivatives forces
every sectional curvature of a plane containing a nonzero gradient to
equal -1/3. CH has different holomorphic and totally real sectional
curvatures at every point, so the gradient vanishes and then f=0. The
uniform nearby inverse in Delay's proposition permits the small inner
metric as reference as well.

Choose mu sufficiently small that K\_\{F\_mu\} is close to CH on the
fixed annulus, then choose a sufficiently small. Both metrics have
scalar -4, so the glued metric has scalar -4 everywhere. The correction
tends to zero in C2 with the annular data; the annulus retains
f\textgreater0, since CH has strictly positive f there. Consequently the
completed cell has

\[
 s=-4,\qquad N_C=1-\delta<1,\qquad J_C<\infty. \tag{4.3}
\]

We can also choose \(J_C>8\pi^2\). At a=0 the unchanged inner core
already has cost \(8\pi^2+8\pi^2\mu+O(\mu^2)\) when truncated at a fixed
radius; the remainder contributes nonnegative cost. Choose a smaller if
necessary and use continuous dependence. Parameters are now fixed, not
sent to zero together with h.

\hypertarget{adding-handles-while-retaining-the-negative-mass-loss}{%
\subsubsection{4.2 Adding handles while retaining the negative-mass
loss}\label{adding-handles-while-retaining-the-negative-mass-loss}}

Insert k\_j copies of (4.3) into disjoint ball lifts in X\_j. This
produces a closed scalar-minus-four metric B\_j on
\(X_j\#k_j\overline{\mathbb{CP}}^2\), equal to CH away from the cells,
and \(N(B_j)=k_j(1-\delta)\). There are enough disjoint lifts because
\(k_j/A_j\to0\).

For each fixed j add the required ell\_j handles by the classical
shrinking scalar-flat conformally-flat neck construction. All necks are
placed in the unchanged CH region. Denote the resulting metric by
g\_\{j,e\}. The construction, with all transition regions included, has
the following fixed-j limits:

\[
 \begin{aligned}
 &g_{j,e}\longrightarrow B_j\quad\text{smoothly off the finitely many mouths},\\
 &\operatorname{Vol}(g_{j,e})\longrightarrow\operatorname{Vol}(B_j),
 \quad\int s_{g_{j,e}}^2\longrightarrow16\operatorname{Vol}(B_j),\\
 &R(g_{j,e})\longrightarrow R(B_j)+32\pi^2\ell_j.
 \end{aligned} \tag{4.4}
\]

Here the scalar curvature can be bounded independently of e for each
fixed j. A scaled scalar-flat neck has zero scalar curvature. Cut its
Euclidean-end error at radius equal to the square root of the neck
scale: its second-derivative error is O(1), the CH Taylor error also
gives O(1), and the transition volume tends to zero. The inserted
self-dual Weyl energy tends to zero. The last limit in (4.4) follows
from the relative identity I54 (1.1), with handle weight four. These are
exactly the classical shrinking handles already used in I44, not a new
scalar-curvature gluing theorem for handles.

Take the negative Yamabe representative
\(\widehat g_{j,e}=u_{j,e}^2g_{j,e}\) with scalar -4. Such a
representative exists by the same all-metric actual positive supply as
before; no new monopole solution on a manifold with boundary is being
assumed. For fixed j the scalar lower bound and maximum principle bound
u uniformly from above. Interior elliptic estimates give subsequential
smooth convergence away from the mouths. The limit extends across those
points: it is a bounded solution of a semilinear elliptic equation on
the smooth B\_j.

The limit is not zero. The actual scalar bound gives
\(16\operatorname{Vol}(\widehat g_{j,e})\ge32\pi^2A_j\), whereas the
uniform upper bound on u makes the volume of the disappearing necks
negligible. The strong maximum principle makes the limiting solution
positive. Uniqueness of the scalar-minus-four representative in a
negative conformal class then forces this limit to be one. Therefore

\[
 u_{j,e}\longrightarrow1\text{ locally smoothly off the mouths},\qquad
 S(g_{j,e})-S(\widehat g_{j,e})\longrightarrow0. \tag{4.5}
\]

The exact conformal identities of I41 now imply
\(R(\widehat g_{j,e})\to R(B_j)+32\pi^2\ell_j\). Subtract the convergent
energies on the complement of small fixed balls, then let those ball
radii shrink. The total remaining P-energy tends to \(32\pi^2\ell_j\).
Since \(8\pi^2N\le P\) on every measurable region,

\[
 \limsup_{e\to0}N(\widehat g_{j,e})
 \le k_j(1-\delta)+4\ell_j. \tag{4.6}
\]

Choose a diagonal e\_j so that all fixed-j errors in (4.4)--(4.6) are
o(Q\_j). The topology is now exactly the original balanced M\_j, with
its inherited actual conformal/weight supply and class-square budget.
Since \(k_j/Q_j\to2/9\), the resulting metrics satisfy

\[
 \boxed{s_{g_j}=-4,\qquad \mathcal E(g_j)=O(h_j),\qquad
 \limsup\frac{N(g_j)}{Q_j}\le1-\frac{2\delta}{9}<1.} \tag{4.7}
\]

Thus Proposition 1's bulk selection does not follow even from the
constant-scalar competitive inputs on the original balanced family. This
is a failure of that stronger proposed intermediate statement, not a
failure of I55 with its hypotheses.

For completeness, the same fixed-j energy limits give

\[
 \frac{\mathcal E(g_j)}{h_j}
 \longrightarrow\frac{28}{9}+\frac{J_C}{9\pi^2}>4. \tag{4.8}
\]

The actual sequence is not asymptotically minimizing: the inherited
upper bound has coefficient at most four. It therefore does not answer
the unrestricted coefficient problem or refute selection for sequences
whose excess over the actual infimum is o(h).

\hypertarget{the-compensation-that-a-sharp-argument-must-retain}{%
\subsection{5. The compensation that a sharp argument must
retain}\label{the-compensation-that-a-sharp-argument-must-retain}}

The sign partition gives an exact nonnegative remainder

\[
 Z(g)=\int_{\{f\ge0\}}|\operatorname{Ric}^{\circ}|^2
       +\int_{\{f<0\}}\left(\frac{s^2}{12}+4|W^+|^2\right),
 \qquad P=8\pi^2N+Z. \tag{5.1}
\]

Consequently

\[
 \mathcal E=\frac{32\pi^2N+4Z+L_s}{R},\qquad
 \frac{\mathcal E}{h}
 =\frac{4N}{Q}+\frac{Z}{2\pi^2Q}+\frac{L_s}{8\pi^2Q}+o(1)
 \tag{5.2}
\]

on every competitive sequence. This is an identity, not another
independent coercivity theorem. It explains why losing negative mass
need not lower the defect. The scalar-minus-four examples above lose
mass and pay a larger nonnegative Z remainder.

For the actual variational target, I41 permits replacing near-minimizers
by their negative Yamabe representatives. In this reduced class the
coefficient-four lower bound is equivalent to the still unproved
compensated inequality

\[
 Z(g_j)\ge8\pi^2\bigl(Q_j-N(g_j)\bigr)-o(Q_j).
 \tag{5.3}
\]

When N already exceeds Q the inequality is automatic; otherwise its
content is the payment by the rest of the actual curvature. Equation
(4.7) shows why one cannot simply delete this case. The genuinely
remaining alternatives are to prove (5.3) from actual geometry or the
actual field equations, to supply the stronger selection only for true
near-minimizers, or to construct a complete cell/sequence whose total
first-order cost is below four. No such sequence is constructed here.
The unrestricted weak zero-Weyl endpoint is also unchanged by these
arguments.

\hypertarget{a-necessary-degeneration-property-of-actual-near-minimizers}{%
\subsubsection{5.1 A necessary degeneration property of actual
near-minimizers}\label{a-necessary-degeneration-property-of-actual-near-minimizers}}

There is a variational consequence that goes beyond relabelling (5.3).
Fix a nontrivial marked cell C of I54, with q\textgreater0 and a fixed
open CH collar. Let K be a compact family in the C4 topology consisting
of smooth positive metrics on C, all with that fixed collar. Allow a
finite union of such families on finitely many fixed cell types.

\textbf{Proposition 2.} Let g\_j be scalar-minus-four metrics on the
original balanced M\_j with

\[
 \mathcal E(g_j)-\alpha_j=o(h_j),\qquad
 \alpha_j=\inf_{g\text{ on }M_j}\mathcal E(g).
 \tag{5.4}
\]

If g\_j contains n\_j disjoint marked full-collar cells whose pullback
metrics lie in K, then \(n_j=o(Q_j)\). No hypothesis that a smooth
global minimizer exists is used. More strongly, there is a fixed C4
neighborhood of K with the same conclusion. The actual cells in this
neighborhood need not themselves have exact CH collars; all variations
are supported away from their boundaries.

\emph{Proof.} The inherited upper bound on alpha\_j and (5.4) first give
\(\mathcal E(g_j)=O(h_j)\). Every cell in K has positive J: zero would
make it Einstein, contradicting I54's topology and open-collar argument.
Compactness and continuity give \(J\ge j_0>0\). Since the initial
metrics have s=-4, their cell J equals P. Equation (1.3) therefore gives
n\_j=O(Q\_j).

I54, Proposition 2 proves that every smooth nontrivial full-CH-collar
cell has a compactly supported smooth metric variation strictly
decreasing J. For each metric in K choose one sufficiently small finite
variation. Continuity in the initial metric makes that same variation
decrease J by a fixed positive amount on a neighborhood. A finite
subcover of K therefore supplies a uniform number d\textgreater0 and
finitely many variations such that each cell can be changed with
\(\Delta J\le-d\). Their changes in volume, scalar integral and Ricci
energy are uniformly bounded. All modified metrics remain positive and
keep the complete collar.

Apply those changes simultaneously on the n\_j disjoint cells. Put
\(v=\Delta V\), \(i=\Delta\int s\), \(r=\Delta R\). Their absolute
values are O(n\_j), and

\[
 \Delta J=r+2i+4v\le-d n_j.
\]

The actual identity \(\mathcal E=4-I^2/(VR)\), expanded at \(I=-4V\),
\(R=4V+P\), gives

\[
 \mathcal E(g'_j)-\mathcal E(g_j)
 =\frac{4\Delta J}{R}
  +O\left(\frac{h_jn_j}{A_j}+\frac{n_j^2}{A_j^2}\right).
 \tag{5.5}
\]

For example, its linear remainder is exactly \(-4Pr/R^2\); the remaining
terms are quadratic in the finite moment changes. Bounds (1.3), scalar
-4 and \(n_j/A_j=O(h_j)\) justify the uniform expansion. Consequently,
for some d'\textgreater0 and all large j with n\_j\textgreater0,

\[
 \mathcal E(g'_j)\le\mathcal E(g_j)-d' n_j/A_j.
\]

Although the varied metric need not retain constant scalar curvature, it
is an allowed competitor for alpha\_j. Combining this inequality with
(5.4) proves n\_j=o(Q\_j). QED.

The neighborhood assertion follows from the same finite subcover:
positivity of J, admissibility of the finitely many supported
variations, the strict J decrease and their moment bounds all persist in
a sufficiently small C4 neighborhood of K. The actual metric is
untouched at its boundary, so no new matching argument is needed.

Thus the fixed smooth cells responsible for (4.7) cannot persist with
positive Q-density in a true near-minimizer. A positive-density cell
description of such a sequence, if one exists, must leave every fixed
smooth compact family, vary its collar/scale, or change its geometric
decomposition. The proposition does not extract any cell description,
nor give a uniform descent rate for degenerating cells. That rate,
measured against the missing characteristic mass rather than merely the
number of cells, is a remaining route to (5.3).

\hypertarget{verification-and-next-use}{%
\subsection{6. Verification and next
use}\label{verification-and-next-use}}

\texttt{tools/check\_signed\_characteristic\_mass.py} checks (3.3), the
recovery primitive and its exact negative mass, the minimum's quadratic
coefficient, and the Kähler/Riemann--Roch normalizations. Its actual
complete-cell integration at mu=0.01, outer radius 2.5 and epsilon=0.01
gives N=0.999973955691314 and J/(2 pi squared)=4.077939349922743. The
numerical receipt illustrates Section 3, not the constant-scalar
existence theorem, which uses (4.1), localized scalar deformation and
the fixed-j limiting proof.

The two primary external inputs are Demailly's weak holomorphic Morse
inequality and Delay's localized scalar gluing theorem. The SW supply,
shrinking handles, conformal energy identities, and cover dilution are
inherited with their original assumptions. No full concentration
decomposition, geometric compactness, or universal near-minimizer
classification has been inferred from these inputs.

\end{document}
